\documentclass[11pt]{article}
\usepackage[a4paper,margin=22mm]{geometry}
\usepackage{fontspec}
\usepackage{babel}
\babelprovide[main,import]{hindi}
\babelprovide[import]{english}
\ifdefined\XeTeXgenerateactualtext
  \XeTeXgenerateactualtext=1
\fi
\setmainfont[
  Path=../../../../05_translation/openlogic_hi_9620cc7/locale/hi/fonts/,
  BoldFont=NotoSerifDevanagari-wdth-wght.ttf,
  BoldFeatures={FakeBold=2},
  ItalicFont=NotoSerifDevanagari-wdth-wght.ttf,
  ItalicFeatures={FakeSlant=0.2},
  Script=Devanagari,
  Language=Hindi
]{NotoSerifDevanagari-wdth-wght.ttf}
\setsansfont[
  Path=../../../../05_translation/openlogic_hi_9620cc7/locale/hi/fonts/,
  BoldFont=NotoSerifDevanagari-wdth-wght.ttf,
  BoldFeatures={FakeBold=2},
  Script=Devanagari,
  Language=Hindi
]{NotoSerifDevanagari-wdth-wght.ttf}
\usepackage{xcolor}
\usepackage{tabularx}
\usepackage{booktabs}
\usepackage{enumitem}
\usepackage{hyperref}
\definecolor{ink}{HTML}{17324D}
\definecolor{accent}{HTML}{B54D32}
\definecolor{soft}{HTML}{EFF3F6}
\hypersetup{
  unicode=true,
  pdflang={hi-IN},
  pdftitle={Open Logic Project — Hindi Working Reader (158 of 722 source files)},
  pdfauthor={Open Logic Project contributors; Hindi translation prepared at the direction of Floris},
  pdfsubject={A compiled working reader of the accepted Hindi translation},
  pdfkeywords={Hindi, Devanagari, logic, open textbook, translation},
  colorlinks=true,
  linkcolor=ink,
  urlcolor=accent
}
\pagestyle{empty}
\setlength{\parindent}{0pt}
\setlength{\parskip}{0.65em}
\newcommand{\ruleline}{\textcolor{accent}{\rule{\linewidth}{1.4pt}}}
\newcommand{\module}[4]{%
  #1 & \textbf{#2}\\[-1pt]
  & {\small\foreignlanguage{english}{#3}} & #4\\[5pt]
}

\begin{document}
\pagecolor{soft}
\vspace*{16mm}
{\color{accent}\sffamily\bfseries\large हिंदी • hi-Deva-IN}\par
\vspace{7mm}
{\color{ink}\sffamily\bfseries\fontsize{30}{38}\selectfont
मुक्त तर्क परियोजना}\par
\vspace{2mm}
{\color{ink}\sffamily\fontsize{20}{26}\selectfont
हिंदी कार्यशील पाठक}\par
\vspace{6mm}
{\color{ink}\sffamily\bfseries\Large
\foreignlanguage{english}{Open Logic Project — Hindi Working Reader}}\par
\vspace{14mm}
\ruleline
\vspace{5mm}
{\sffamily\bfseries\fontsize{22}{28}\selectfont
158 / 722 स्रोत फ़ाइलें}\par
{\sffamily\large 59,955 स्रोत-शब्दों का स्वीकृत अनुवाद}\par
\vfill
\begin{minipage}{0.84\linewidth}
यह एक वास्तविक, संकलित पाठक है—केवल प्रगति-सूचना नहीं। इसमें नीचे सूचीबद्ध
158 स्रोत फ़ाइलों का हिंदी पाठ है। संपूर्ण 722-फ़ाइल परियोजना का अनुवाद अभी
जारी है।
\end{minipage}
\vfill
{\small\sffamily
स्थिर स्रोत संशोधन: \texttt{9620cc7}\quad•\quad 15 अगस्त 2026\par
क्रिएटिव कॉमन्स नामोल्लेख 4.0 अंतरराष्ट्रीय (CC BY 4.0)}
\clearpage
\pagecolor{white}

{\color{ink}\sffamily\bfseries\LARGE इस पाठक के बारे में}\par
\ruleline

यह संस्करण मुक्त तर्क परियोजना के स्थिर अंग्रेज़ी स्रोत संशोधन
\texttt{9620cc73f9c8e0ad003c514a5d3748f29611c4c0} (वृक्ष
\texttt{f67757bb9305b173634082ab4cefd5601a707a34}) पर आधारित है। यहाँ शामिल
158 सामग्री-फ़ाइलों का पाठ देवनागरी हिंदी में अनूदित है; उनका मापा हुआ स्रोत
पाठ 59,955 शब्द है।

\textbf{क्या पूरा है।} सूचीबद्ध फ़ाइलों के पाठ का अनुवाद, प्रत्येक स्वीकृत
खंड का XeLaTeX संकलन, खोज/प्रतिलिपि जाँच, तथा प्रस्तुत पृष्ठों की दृश्य जाँच
की गई है। गणितीय सूत्र, LaTeX आदेश, स्थिर पहचानकर्ता, लेबल और उद्धरण-संरचना
को स्रोत के अनुरूप रखा गया है।

\textbf{क्या अभी पूरा नहीं है।} यह संपूर्ण मुक्त तर्क परियोजना का हिंदी
संस्करण नहीं है। कुल 722 सामग्री-फ़ाइलों में से 158 यहाँ सम्मिलित हैं। शेष
अनुवाद क्रमशः जोड़े जाएँगे। इसे मानव-समीक्षित आलोचनात्मक संस्करण या PDF/UA
प्रमाणन न माना जाए।

\textbf{श्रेय और परिवर्तन।} मूल सामग्री मुक्त तर्क परियोजना के लेखकों और
योगदानकर्ताओं की है। हिंदी अनुवाद Floris के निर्देशन में OpenAI 5.6 Sol
(Ultra mode) द्वारा तैयार किया गया। यह रूपांतरण मूल से परिवर्तन है। स्रोत:
\href{https://github.com/OpenLogicProject/OpenLogic}{GitHub पर मूल परियोजना}।

मूल और अनूदित सामग्री
\href{https://creativecommons.org/licenses/by/4.0/}{CC BY 4.0} के अधीन दी
गई है। पाठक PDF को पहले रखा गया है; संपादनयोग्य स्रोत, फ़ाइल-स्तरीय हैश,
निर्णय-पुस्तिका और QA प्रमाण अलग डाउनलोडों में उपलब्ध हैं।

\vfill
{\color{ink}\sffamily\bfseries\large उपयोग के लिए संक्षिप्त सूचना}\par
यह कार्यशील संस्करण पढ़ने और प्रतिक्रिया देने के लिए प्रकाशित है। किसी
विशिष्ट अंश का उल्लेख करते समय इस संस्करण की DOI और पृष्ठ संख्या दें।
\clearpage

{\color{ink}\sffamily\bfseries\LARGE विषय-सूची}\par
{\small भौतिक PDF पृष्ठ; कोष्ठक में सम्मिलित स्रोत फ़ाइलें।}\par
\ruleline
\vspace{3mm}
\renewcommand{\arraystretch}{1.05}
\begin{tabularx}{\linewidth}{@{}r X r@{}}
\textbf{क्रम} & \textbf{भाग} & \textbf{पृष्ठ}\\
\midrule
\module{1}{समुच्चय}{Sets • 7 files}{5–13}
\module{2}{संबंध}{Relations • 10 files}{14–23}
\module{3}{फलन}{Functions • 7 files}{24–32}
\module{4}{समाकृतिक फलन — पूरक पाठ}{Isomorphic functions • 1 file}{33}
\module{5}{समुच्चयों का आकार}{Size of sets • 15 files}{34–54}
\module{6}{अंकगणितीकरण}{Arithmetization • 8 files}{55–68}
\module{7}{अनंत समुच्चय}{Infinite sets • 6 files}{69–75}
\module{8}{प्रतिज्ञप्ति तर्क: वाक्यविन्यास और अर्थविज्ञान}{Propositional logic: syntax and semantics • 7 files}{76–84}
\module{9}{प्रथम-क्रम तर्क का परिचय}{Introduction to first-order logic • 10 files}{85–92}
\bottomrule
\end{tabularx}
\vfill
{\small “पूरक पाठ” उस स्वीकृत स्रोत-फ़ाइल को दर्शाता है जिसे वर्तमान अध्याय
ड्राइवर सीधे आयात नहीं करता; उसे पाठक से गायब करने के बजाय अलग संकलित करके
ठीक स्थान पर जोड़ा गया है।}
\clearpage

{\color{ink}\sffamily\bfseries\LARGE विषय-सूची — जारी}\par
\ruleline
\vspace{3mm}
\begin{tabularx}{\linewidth}{@{}r X r@{}}
\textbf{क्रम} & \textbf{भाग} & \textbf{पृष्ठ}\\
\midrule
\module{10}{प्रथम-क्रम तर्क: वाक्यविन्यास}{First-order logic: syntax • 10 files}{93–106}
\module{11}{निष्पादन-तंत्र}{Proof systems • 6 files}{107–111}
\module{12}{अनुक्रम-कलन}{Sequent calculus • 15 files}{112–132}
\module{13}{स्वाभाविक निगमन}{Natural deduction • 14 files}{133–155}
\module{14}{सत्य-वृक्ष}{Tableaux • 14 files}{156–178}
\module{15}{अभिगृहीतीय निष्पादन}{Axiomatic deduction • 14 files}{179–190}
\module{16}{सिद्धता — पूरक पाठ}{Provability • 1 file}{191–194}
\module{17}{पूर्णता प्रमेय}{Completeness theorem • 12 files}{195–209}
\module{18}{अधिकतम संगत समुच्चय — पूरक पाठ}{Maximally consistent sets • 1 file}{210–211}
\bottomrule
\end{tabularx}

\vfill
\begin{minipage}{\linewidth}
\colorbox{soft}{\parbox{0.94\linewidth}{\small
\textbf{कुल:} 158 स्वीकृत स्रोत फ़ाइलें • 207 अनूदित सामग्री-पृष्ठ •
4 परिचयात्मक पृष्ठ • 211 PDF पृष्ठ।}}
\end{minipage}

\end{document}
